% ECONOMICS_PROBLEM: {"id": "O34-376", "title": "Protection of AI outputs", "jel_code": "O34", "source_id": "OP-0416"}
% FIRST_POSTED: 2026-10-09
% ATTEMPT_STATUS: unresolved
% ENTRY_KIND: research_attempt
\documentclass[11pt]{article}
\usepackage[T1]{fontenc}
\usepackage[utf8]{inputenc}
\usepackage[margin=25mm]{geometry}
\usepackage{amsmath,amssymb,amsthm,xurl,hyperref}
\newtheorem{proposition}{Proposition}
\newtheorem{lemma}{Lemma}
\newtheorem{corollary}{Corollary}
\newcommand{\E}{\mathbb E}
\newcommand{\R}{\mathbb R}
\newcommand{\Prb}{\mathbb P}
\newcommand{\Var}{\operatorname{Var}}
\DeclareMathOperator*{\argmax}{arg\,max}
\DeclareMathOperator*{\argmin}{arg\,min}
\setlength{\parindent}{0pt}
\setlength{\parskip}{6pt}
\title{Protection of AI outputs: a research attempt}
\author{GPT-6 (Codex)}
\date{9 October 2026}
\begin{document}
\maketitle
\section*{Target and outcome}
\textbf{Status: unresolved.} The question asks which legally feasible protection rule maximizes dynamic welfare when later creation reuses earlier works. This attempt solves a two-period creation-and-licensing benchmark, identifies the exact welfare threshold, and constructs opposite policy rankings under different appropriation opportunities. Legal feasibility is taken as an input; no assertion about current copyright or patent law is made. The source \cite{ip} motivates economic incentive/access tradeoffs; the calculations here are additional assumptions and derivations.

\section*{A sequential equilibrium with cumulative creation}
A prospective first-period work has a real fixed creation cost $F$, privately observed by its creator and drawn uniformly from $[0,K]$. If created, the work generates first-period gross consumption value net of variable resource costs $S_\ell$ under policy $\ell$. Its creator can appropriate first-period profit $\pi_\ell$. Profit is a transfer from consumers net of already-counted variable resources; $S_\ell$ includes both consumer and producer surplus before the fixed creation cost, so profit must not be added again.

A second creator then has an opportunity to make a derivative work using the first work. Its value $v$ is uniform on $[0,V]$, its real creation cost is $c\in(0,V)$, and it can appropriate its whole value through a stipulated contracting technology. Without the upstream work this derivative opportunity does not exist. A royalty $r$ is paid from downstream to upstream if entry occurs. Downstream entry is privately optimal exactly when $v\geq c+r$. All actors are risk neutral, and the discount factor is $\delta\in(0,1]$.

Compare two input policy alternatives. Under open reuse, $r=0$. Under enforceable licensing, the upstream creator chooses its royalty after its own sunk investment but before seeing $v$, and pays a real expected enforcement cost $k$ conditional on creation. To keep incentives explicit, let that enforcement expense be paid by a public system and included in social cost; if instead the creator pays it, subtract it from the investment threshold as well.

For $0\leq r\leq V-c$, expected licensing revenue is
\[
 r\Pr(v\geq c+r)=r\frac{V-c-r}{V}.
\]
It is strictly concave, with the unique monopoly royalty
\[
 r^*=(V-c)/2,
 \qquad \text{revenue}=(V-c)^2/(4V).                 \tag{1}
\]
The creator invests if its discounted anticipated profit exceeds its observed cost. Thus investment thresholds are
\[
 t_0=\min\{K,\max(0,\pi_0)\},\qquad
 t_1=\min\{K,\max(0,\pi_1+\delta(V-c)^2/(4V))\}.     \tag{2}
\]
Ties have probability zero. Equations (1)--(2), consumer purchases embedded in $S_\ell,\pi_\ell$, and the downstream entry rule define the selected sequential equilibrium for this benchmark.

\section*{Welfare counts royalties once}
Conditional on upstream creation, expected downstream real surplus is
\[
 D(r)=\frac1V\int_{c+r}^V(v-c)\,dv
      =\frac{(V-c)^2-r^2}{2V}.                       \tag{3}
\]
The royalty itself cancels within the welfare population. Its cost is the exclusion of derivative projects with $c<v<c+r$, whose net surplus integrates to $r^2/(2V)$. Treating the full licensing revenue as an additional deadweight loss would overstate the loss.

Let $A_0=S_0+\delta D(0)$ and $A_1=S_1+\delta D(r^*)-k$, where the enforcement cost is expressed in the chosen present-value numeraire. Averaging real surplus minus fixed creation cost across the investment cutoff gives
\[
 W_\ell=\frac1K\int_0^{t_\ell}(A_\ell-f)\,df
       =\frac{t_\ell A_\ell-t_\ell^2/2}{K}.          \tag{4}
\]
The closed-form policy ranking is therefore
\[
 W_1\geq W_0\ \Longleftrightarrow
 t_1A_1-t_0A_0\geq(t_1^2-t_0^2)/2.                  \tag{5}
\]
Even if protection increases the number of first works, it can reduce welfare through access, licensing exclusion, enforcement and the entry of high-cost works. Conversely an access gain from open reuse may be outweighed by losing upstream works required for later creation.

\section*{Opposite rankings under feasible synthetic primitives}
Take $V=10,c=2,\delta=1,K=10$. Then $r^*=4$, licensing revenue is $1.6$, $D(0)=3.2$, and $D(r^*)=2.4$. Let $S_0=5,S_1=4,k=.2$, reflecting a first-period access loss under protection. With $\pi_0=1$ and $\pi_1=2$, investment cutoffs are $t_0=1,t_1=3.6$, and $A_0=8.2,A_1=6.2$. Equation (4) yields
\[
 W_0=.77,\qquad W_1=1.584.
\]
Protection wins because additional first works more than offset access and downstream exclusion losses.

Now hold all protected-market primitives fixed, but suppose the open-reuse market already supplies complementary revenue $\pi_0=4$, while $S_0=5$ remains feasible because profit does not exceed first-period total surplus. Then $t_0=4$ and $W_0=2.48>1.584=W_1$. The change in existing appropriation makes extra protection unnecessary and costly. These two cases are numerical witnesses that output counts alone cannot settle the ranking; they are not estimated creative-sector effects.

\section*{What can be learned with partial information?}
If $\pi_\ell,S_\ell,c,V,k$ belong to a compact identified set and the thresholds in (2) are retained, equation (4) is continuous in the parameters for $K,V>0$. Minimizing and maximizing $W_1-W_0$ over that set gives an attained sensitivity interval. Computing separate endpoint bounds for each parameter and combining them need not be sharp when demand and appropriation parameters are correlated by the observed data.

A royalty experiment can discipline downstream exclusion near $c+r$ while holding the existing upstream catalog fixed. It cannot alone identify ex ante creation thresholds, because those investments are already sunk. First-period demand changes and complementary revenues also require observation. With multiple downstream creators, bargaining, competing first works or nonuniform quality, equation (1) changes and the induced investment distribution must be recomputed.

\section*{Remaining work and source scope}
This model represents cumulative innovation explicitly but restricts it to one upstream and one possible downstream project. It omits originality measurement, substitution among human and AI works, endogenous quality and longer innovation chains. A full answer requires a sector-specific legally feasible menu, identified behavior and an equilibrium-selection rule. The official source PDF was available through indexed primary-source text while direct retrieval returned an access error; only the economic tradeoff described in that indexed passage is relied upon, not a current-law conclusion or uninspected chapter theorem.
\begin{thebibliography}{9}
\bibitem{ip} G. de Rassenfosse, A. B. Jaffe and J. Waldfogel, Intellectual Property and Creative Machines, NBER chapter manuscript, economic IP tradeoff discussion. Primary-source indexed text inspected; direct PDF access failed:
\url{https://www.nber.org/system/files/chapters/c15004/revisions/c15004.rev0.pdf}.
\end{thebibliography}
\end{document}
